\section{引言}

通行的读法是：明代的皇权被一套「家国同构」的伦理语言包住——
皇帝既是政治主权也是天下之父，官僚是父母官，臣之忠与子之孝共用同一套代码。
这套读法靠的是引文与印象，很少被点数。本篇把问题换成可计数的：在这部被读过最多的
明代通俗叙事里，「家」和「继承」到底是靠哪几个字被记住的，又把这些字所指的人
放进了什么样的关系网。语料取自同一部书的结构化投影：1231 位人物、2324 条关系、
563 处地点。我们量的不是这套说法是否符合明代实况，而是这部叙事在写它时
用了哪些词、把这组词的人挂在了什么样的边上。

具体命题是：\textbf{「家」究竟是不是这部叙事的组织原则——词汇侧那点压倒性的高频
（族一次命中 234 次，居全书词表之首），能不能在关系结构上换来一个亲缘主导的子图？} 若两者在关系类别构成上确实分离，那么这一母题就不是后人的印象收拾，而是这部叙事自己组织出来的认知结构。

语料是通俗明史叙事《明朝那些事儿》全七部（全书 156 章，本篇锚点覆盖 41 章）及从同一语料抽取、可独立重建的结构化知识图谱：人物 1231、事件 1106、关系 2324、地点 563。我们量的不是该母题是否属实，而是这部叙事在写它时用了哪些词、把这组词的人放进了什么样的网络。

结果指向一个清楚的形状。四个字就把整个母题的光照住了：族、孝、嗣、继承一抬起来，
剩下的词项只是陪衬，撑不起一次统计检验。网络一侧，以四个核心锚点拉起的 130 人子图含显式边 636 条，四个锚点的自我网络把大半的皇权枢纽连成一块，类别构成因此出现一处清晰位移：
未归类的「其他」一路抬高，而亲缘这一类几乎原地不动。叙事热度与网络度数的 Spearman 相关为 $\rho=0.871$，近乎同一个排布——在这部叙事里，常被写到的人就是连得多的人，两条线没有分岔。均值层面，子图成员平均度数 12.6，全图平均 6.3。这两个均值不可并读：子图是围绕四个超级枢纽搭起来的，12.6 是枢纽邻居的均值，
不是子图内部更连通的证据；全图那个均值里还压着 537 个孤立点。把两者相减之前，
先要承认它量的是选集偏差，不是密度差异。

贡献有三点：
\begin{itemize}
\item 给出一套十二词的继承与亲缘语汇表（另含两个零命中词），逐字命中口径写死、
计数可复算；这是首次对这一母题做词汇侧的硬点数。
\item 把母题从印象搬进关系结构：以四个核心人物的自我网络（130 人、636 条显式边）
报告它与全图基线之间逐类别的位移，而不是只报一个「亲属很多」的判断。
\item 报一个与词汇侧直觉相反、但可复核的结果——亲缘边占比几乎不动（子图 7.5 对基线 8.0），
把「词汇可见性不等于关系组织度」落成两个能互相拆台的小数。
\end{itemize}

边界先说死：本文不主张这套说法符合该领域的既有结论，也不主张它符合任何一人的真实心态。它量的是一部被大量阅读的叙事如何组织关于这一母题的印象。

第 2 节回顾四支相关文献，第 3 节交代语料与口径，第 4 节写方法，第 5 节报告结果，第 6 节讨论效度边界，第 7 节收束。
